\documentclass[11pt]{article}

\usepackage[margin=1in]{geometry}
\usepackage{booktabs}
\usepackage{amsmath}
\usepackage{graphicx}
\usepackage[hidelinks]{hyperref}
\usepackage[numbers,sort&compress]{natbib}

\title{External Paired-Scanner Validation of PathoAlign on Canine Squamous Cell Carcinoma Histopathology}
\author{Larry Chorne}
\date{\today}

\begin{document}
\maketitle

\begin{abstract}
Scanner-specific acquisition effects can confound representation learning in computational pathology, especially when models are evaluated across sites, devices, or staining/scanning pipelines. PathoAlign is a paired-acquisition projection objective designed to preserve biological region identity while reducing scanner-identifiable signal in the biological embedding. Here, we evaluate whether a PathoAlign configuration selected on a human paired-scanner benchmark transfers to an independent public multi-scanner canine cutaneous squamous cell carcinoma dataset. After auditing annotation correspondence and applying deterministic geometry qualification, we retained 805 complete five-scanner regions from 44 biological samples. Using frozen DINOv2-Base features and locked hyperparameters, PathoAlign reduced sample-blocked scanner-probe accuracy by 0.380 absolute while improving cross-scanner pair cosine and preserving same-region retrieval. These results support PathoAlign as a paired-acquisition representation method that transfers beyond the development benchmark, while remaining a research-only validation rather than a clinical or diagnostic claim.
\end{abstract}

\section{Introduction}

Computational pathology representations can encode both tissue morphology and acquisition artifacts. When acquisition signal is correlated with scanner, site, or cohort membership, downstream performance may reflect technical shortcuts rather than robust biological representation. Paired-acquisition datasets provide a direct way to test whether a representation preserves tissue identity while reducing scanner-identifiable variation.

PathoAlign separates a frozen image representation into a biological embedding and an acquisition embedding by combining paired-region contrastive alignment, reconstruction, variance/covariance regularization, scanner-adversarial learning, and scanner-dependence penalties. The present study asks whether the locked PathoAlign configuration transfers to an independent external paired-scanner benchmark.

\section{Dataset and geometry qualification}

We use the public Multi-Scanner Canine Cutaneous Squamous Cell Carcinoma histopathology dataset, which contains 44 biological samples scanned on five scanners: CS2, GT450, NZ20, NZ210, and P1000. The raw release contains 220 TIFF images and 6,215 annotation views corresponding to 1,243 candidate matched regions.

Before modeling, we audited the COCO annotations, SlideRunner SQLite database, TIFF dimensions, and scanner-to-scanner geometry. The P1000 view showed a coherent sample-specific orientation discrepancy rather than random annotation failure. We therefore applied deterministic P1000 orientation normalization and retained only regions for which every scanner view required at most 10\% adaptive-crop padding. This produced a geometry-qualified paired-acquisition subset of 805 complete five-view regions, corresponding to 4,025 image views.

\section{Frozen encoder triage}

We evaluated frozen ResNet50 ImageNet, DINOv2-Base, and Phikon representations on the fold-0 validation split before running PathoAlign. DINOv2-Base was selected as the external substrate because it had the strongest cross-scanner cosine consistency and best worst-case retrieval while still retaining scanner-identifiable signal above chance.

\begin{table}[ht]
\centering
\caption{Frozen encoder triage on fold-0 validation.}
\label{tab:encoder-triage}
\input{tables/encoder_triage.tex}
\end{table}

\section{Locked PathoAlign protocol}

The PathoAlign hyperparameters were locked before the external five-fold test. The validation run used seeds 901--905 on fold 0, fitting on train and evaluating on validation. The final external test used five sample-blocked folds and new seeds 911--915. For each fold, all non-test samples were used for fitting and the held-out fold was projected and evaluated exactly once. The two variants were a paired-consistency reference projection and the locked PathoAlign dependence-penalty configuration, denoted PathoAlign dep20.

\section{Results}

Across five folds and five seeds, PathoAlign reduced scanner identifiability while preserving same-region retrieval. Descriptive run means are shown in Table~\ref{tab:run-means}. Sample-blocked paired contrasts over 44 biological samples are shown in Table~\ref{tab:sample-contrasts}.

\begin{table}[ht]
\centering
\caption{Locked five-fold descriptive run means.}
\label{tab:run-means}
\input{tables/descriptive_run_means.tex}
\end{table}

\begin{table}[ht]
\centering
\caption{Sample-blocked PathoAlign minus paired-reference contrasts over 44 biological samples.}
\label{tab:sample-contrasts}
\input{tables/sample_blocked_contrasts.tex}
\end{table}

The main scanner-probe contrast was $-0.380$ absolute with a bootstrap 95\% confidence interval of $[-0.400, -0.361]$ and Monte Carlo sign-flip $p=0.000004$. All 44 sample blocks favored PathoAlign for scanner-probe reduction. Pair-cosine average and worst-pair cosine also improved for all 44 samples. Retrieval top-1 average and worst-pair retrieval remained non-inferior within the predefined 0.02 margin.

The acquisition factor retained scanner information, with mean acquisition scanner-probe accuracy of 0.865, while acquisition-space tissue retrieval was low at 0.181. The mean biological--acquisition cross-covariance RMS was 0.0898. This pattern is consistent with scanner signal being preferentially represented in the acquisition branch rather than the biological embedding.

\section{Discussion}

This external study supports the transferability of PathoAlign's paired-acquisition objective. The method reduced scanner-identifiable signal by a large absolute margin on a dataset not used for hyperparameter selection, while improving cross-scanner cosine consistency and preserving same-region retrieval. The result is important because it demonstrates that the effect is not confined to the original human paired-scanner benchmark.

The study has several limitations. The dataset is canine SCC rather than human pathology, and the experiment evaluates representation identifiability rather than clinical performance. The geometry-qualified subset intentionally excludes regions that required substantial padding or unreliable crop support. Finally, PathoAlign is evaluated here as a projection on frozen features, not as an end-to-end diagnostic model.

\section{Conclusion}

A PathoAlign configuration locked from prior development transferred to an independent external paired-scanner canine SCC benchmark. On 44 sample-blocked biological units, PathoAlign reduced scanner-probe accuracy by approximately 0.38 absolute, improved cross-scanner cosine consistency, and preserved same-region retrieval. This provides external paired-acquisition evidence for PathoAlign as a research method for scanner-aware representation identifiability in computational pathology.

\bibliographystyle{plainnat}
\bibliography{references}

\end{document}
